% ID: SEL-D0084-OB-E01
% Tanda S084.
Sean $x$ e $y$ los números de estanterías grandes y pequeñas. Las restricciones de madera, producción mínima y proporción entre tipos son
\[
 4x+3y\le60,\qquad x\ge3,\qquad y\ge2x.
\]
Estas condiciones implican $y\ge6$ y delimitan un triángulo cerrado y acotado. Sus vértices son $(3,6)$, $(3,16)$ y $(6,12)$. El último surge de $4x+3(2x)=60$.

Maximizamos $Z=60x+40y$, en euros. En los vértices,
\[
 Z(3,6)=420,\qquad Z(3,16)=820,\qquad Z(6,12)=840.
\]
Como la función es lineal, el máximo es $\boxed{840}$ euros y se obtiene con 6 estanterías grandes y 12 pequeñas. Son cantidades enteras, consumen $4\cdot6+3\cdot12=60$ metros cuadrados de tablero y cumplen la proporción mínima de pequeñas.
\begin{center}
\setlength{\unitlength}{1mm}
\begin{picture}(106,77)(-10,-10)
\put(-4,0){\vector(1,0){96}}\put(0,-5){\vector(0,1){66}}
\put(94,0){$x$}\put(1,65){$y$}
\qbezier(30,18)(30,33)(30,48)
\qbezier(30,48)(45,42)(60,36)
\qbezier(60,36)(45,27)(30,18)
\put(13,13){\scriptsize $(3,6)$}
\put(15,51){\scriptsize $(3,16)$}
\put(61,36){\scriptsize $(6,12)$}
\end{picture}
\end{center}
